Cieľom tejto diplomovej práce bolo navrhnúť a implementovať inteligentného
finančného asistenta FinanceGPT, integrovaného do existujúcej webovej aplikácie na
správu faktúr. Asistent umožňuje používateľom komunikovať s finančnými dátami
prostredníctvom prirodzeného jazyka, pričom výstupom sú textové odpovede,
numerické agregácie aj vizuálne reprezentácie dát.

V teoretickej a analytickej časti sme sa zamerali na súčasné možnosti
spracovania prirodzeného jazyka s využitím LLM. Ako spôsob integrácie sme zvolili
architektúru Tool Calling, ktorá na rozdiel od klasických klasifikátorov
ponecháva rozhodovanie o použití nástrojov priamo jazykovému modelu. Tento prístup
odbremeňuje systém od potreby definovať všetky používateľské scenáre vopred.

Architektúra systému rešpektuje konvencie frameworku Ruby on Rails. Konverzačné
rozhranie sme implementovali pomocou technológie Hotwire (Turbo a Stimulus), čím sme
sa vyhli potrebe budovať SPA architektúru. Generovanie SQL dopytov a výpočet
agregácií zostali izolované na strane backendu, pričom asistent slúži ako prekladač
medzi prirodzeným jazykom a relačnou databázou.

V implementačnej fáze bol vytvorený funkčný prototyp. Asistent dokáže prijímať
textový vstup, dynamicky extrahovať kontext z databázovej schémy a pri vizuálnych
požiadavkách generovať štruktúrovaný JSON pre klientské vykreslenie grafov pomocou
Chart.js.

Testovanie na evaluačnej sade 30 dopytov ukázalo, že model GPT-4o-mini dosiahol
úspešnosť 93,3 \% a model Claude 3.5 Haiku 90,0 \%. Najväčšie limity sa neprejavili
v syntaktickom generovaní SQL, ale v sémantickom mapovaní doménových pojmov na
databázovú schému a v kolíziách medzi nástrojmi pri nejednoznačných dopytoch. Výsledky
zároveň potvrdili, že guardraily a prompt engineering majú väčší vplyv na
úspešnosť systému než samotná voľba modelu.

V priebehu testovania ma prekvapilo, že modely dokázali volať viacero nástrojov v
jednom dopyte aj bez špecifického systémového promptu, rovnako ako aj spôsob, akým
dokázali zdôvodňovať odpovede pri neúplnom kontexte.

Z pohľadu ďalšieho vývoja sa ako najdôležitejšia ukazuje oblasť bezpečnosti.
Aktuálna verzia obsahuje validáciu SQL operácií a tenant filtrovanie, avšak pre produkčné
nasadenie by bolo vhodné doplniť pokročilejšie auditovanie databázových operácií a
validáciu generovaných dotazov. Prínosné by bolo aj rozšírenie evaluačného
datasetu a testovanie systému na väčšom objeme reálnych dát.

Primárny cieľ práce bol splnený. Navrhnutý a implementovaný bol prístup, ktorý
posúva interakciu používateľa s fakturačným systémom od tradičných systémov
smerom ku konverzačnému rozhraniu. Výsledný prototyp predstavuje základ pre ďalší
vývoj využívajúci LLM v oblasti finančných aplikácií.